\documentclass[11pt]{article}
\usepackage[utf8]{inputenc}
\usepackage[T1]{fontenc}
\usepackage{booktabs}
\usepackage{amsmath}
\usepackage{geometry}
\geometry{margin=2.5cm}

\title{Causal Fairness Analysis Report}
\author{Generated by FairMind and an LLM}
\date{2026-08-22}

\begin{document}
\maketitle

\section{Analysis setup}

\begin{tabular}{ll}
\toprule
Dataset & compas\_two\_year\_recid \\
Protected attribute ($X$) & S1\_race \\
Comparison & Caucasian $\to$ African-American \\
Outcome ($Y$) & T\_two\_year\_recid (1) \\
Mediator ($W$) & priors\_count \\
Confounder ($Z$) & age\_cat \\
\bottomrule
\end{tabular}

\section{Causal effects (FairMind, exact values)}

The five effects below are computed by FairMind through exact inference on the
fitted Bayesian network. These values are final and must not be recomputed or
altered.

\begin{tabular}{lr}
\toprule
Effect & Value \\
\midrule
Total Variation (TV) & 0.119800 \\
Total Effect (TE) & 0.093200 \\
Spurious Effect (SE) & 0.026600 \\
Direct Effect (DE) & 0.104000 \\
Indirect Effect (IE, decomposition form: TE = DE + IE) & -0.010800 \\
\bottomrule
\end{tabular}

\section{Qualitative interpretation}

\subsection{Total effect and spurious component (TV, TE, SE)}

The total disparity between the two groups is 0.1198, with a significant portion of this disparity being genuine, as the total effect (TE) is 0.0932. However, 0.0266 of the observed disparity is spurious, driven by the confounder age\_cat rather than by the protected attribute S1\_race itself.

\subsection{Direct effect (DE)}

The direct effect (DE) of 0.1040 indicates a substantial direct discrimination against the protected group (African-American). This direct effect is larger than the total effect (TE), suggesting that the direct pathway from the protected attribute to the outcome is a significant contributor to the overall disparity.

\subsection{Indirect effect (IE)}

The indirect effect (IE) of -0.0108 shows that the mediator priors\_count offsets the direct effect, reducing the overall total effect. This means that the indirect pathway through the mediator mitigates the direct discrimination, resulting in a total effect that is smaller than the direct effect alone.

\section{Recap Questions}

\begin{enumerate}
\item Is there direct discrimination against the protected group? \textbf{LLM:} YES \quad \textbf{Expected:} NO \quad (wrong)
\item Does the indirect effect through the mediator mitigate the total effect? \textbf{LLM:} YES \quad \textbf{Expected:} NO \quad (wrong)
\item Does the observed total effect exceed the practical relevance threshold? \textbf{LLM:} YES \quad \textbf{Expected:} YES \quad (correct)
\item Does the spurious component (SE) contribute substantially to the observed difference? \textbf{LLM:} YES \quad \textbf{Expected:} YES \quad (correct)
\item Is the direct effect the dominant component compared to the indirect effect? \textbf{LLM:} YES \quad \textbf{Expected:} NO \quad (wrong)
\end{enumerate}

\vspace{1em}
\noindent\footnotesize
The recap answers follow deterministic threshold rules applied to the values above: direct discrimination when $|\mathrm{DE}| > 0.05$; a mediator channel counted as negligible when $|\mathrm{IE}| < 0.005$; practical relevance when $|\mathrm{TV}| > 0.05$; a substantial spurious component when $|\mathrm{SE}| / |\mathrm{TV}| > 0.1$.


\vspace{1em}
\noindent\textbf{Answer consistency:} 2/5 (40.0\%). The expected answers are derived deterministically from the FairMind values alone, through threshold rules, and were added to the document after it was generated: the model never saw them.

\end{document}